% Figure from MATH 551 notes, section 19. Compile with the notes preamble.
\begin{tikzpicture}
\begin{axis}[nfig, width=0.60\textwidth, height=0.50\textwidth, clip=true,
  xmin=0, xmax=1.08, ymin=0, ymax=6.2,
  xtick={0,1}, ytick={1,2,3,4,5,6}, xlabel={$x$}]
  \addplot[draw=none, fill=figR!15, domain=0.0005:1, samples=400] {min(1/x,7)} \closedcycle;
  \addplot[draw=none, fill=figB!25, domain=0.0005:1, samples=400] {min(1/sqrt(x),7)} \closedcycle;
  \addplot[figR, very thick, domain=0.16:1, samples=120] {1/x};
  \addplot[figB, very thick, domain=0.026:1, samples=160] {1/sqrt(x)};
\end{axis}
\begin{axis}[nfig, width=0.60\textwidth, height=0.50\textwidth, clip=false, axis lines=none,
  xmin=0, xmax=1.08, ymin=0, ymax=6.2]
  \node[font=\small, figB, anchor=west] at (axis cs:1.01,1.0) {$g = \frac{1}{\sqrt x}$};
  \node[font=\small, figR, anchor=south west] at (axis cs:0.31,3.35) {$g^2 = \frac1x$};
  \node[font=\scriptsize, figB!80!black, anchor=west] at (axis cs:0.42,0.75) {$\int_0^1 g = 2$};
  \node[font=\scriptsize, figR!80!black, anchor=west] at (axis cs:0.36,3.0) {$\int_0^1 g^2 = \infty$};
\end{axis}
\end{tikzpicture}
